\documentclass[11pt]{article}
\usepackage[margin=1.1in]{geometry}
\usepackage{amsmath,amssymb,amsthm}
\usepackage{booktabs}
\usepackage{enumitem}
\usepackage[round]{natbib}
\setcitestyle{aysep={}}
\usepackage{xcolor}
\definecolor{linknavy}{RGB}{25,55,120}
\usepackage[colorlinks=true,linkcolor=linknavy,citecolor=linknavy,urlcolor=linknavy,
pdftitle={Coalitional monotonicity and the core: minimal impossibility, an
exact-games escape, and the boundary of surplus-distributor solutions},
pdfauthor={Maxwell Flanagan}]{hyperref}

\newtheorem{theorem}{Theorem}[section]
\newtheorem{lemma}[theorem]{Lemma}
\newtheorem{proposition}[theorem]{Proposition}
\newtheorem{corollary}[theorem]{Corollary}
\newtheorem{conjecture}[theorem]{Conjecture}
\newtheorem{definition}[theorem]{Definition}
\newtheorem{example}[theorem]{Example}
\newtheorem{remark}[theorem]{Remark}
\newtheorem{observation}[theorem]{Observation}

\newcommand{\R}{\mathbb{R}}
\newcommand{\core}{\operatorname{Core}}
\newcommand{\SD}{\operatorname{SD}}

\title{Coalitional monotonicity and the core:\\
minimal impossibility, an exact-games escape,\\
and the boundary of surplus-distributor solutions}

\author{Maxwell Flanagan\\[2pt]
\normalsize Corresponding author: \texttt{maxwell.flanagan@gmail.com}%
\thanks{A machine-readable certificate supplement with a standalone
one-command verifier accompanies this paper as Online Resource~1; see
Section~\ref{sec:methods} for the verification methodology and the
disclosure of AI assistance.}}

\date{July 2026}

\begin{document}
\maketitle

\begin{abstract}
We study coalitional monotonicity of single-valued core selections in
transferable-utility (TU) games. First, we pin down the minimal size of an
impossibility certificate for the classical incompatibility of core
selection and coalitional monotonicity: three-game certificates (a base
game and two single-coalition increases) appear already in the textbook
proof of Peleg and Sudh\"olter, and we record the complementary elementary
fact that two-game families are always feasible---so three is exactly
minimal---together with a self-contained three-game certificate admitting a
six-line proof. Second, we map
the monotonicity boundary of the Arin--Katsev surplus-distributor (SD)
prenucleolus: coalitional monotonicity fails already on \emph{exact} games
($n=4$, a minimized witness), the SD-prekernel is multi-valued on exact
games, population monotonicity fails even on convex games ($n=3$, minimal),
and on permutation games the SD solution amplifies the nucleolus's failure
of pairwise monotonicity, whereas on assignment games extensive searches
support monotonicity of both solutions. Third, we study whether the
impossibility descends to exact games. For $n\le4$ it does not: by known
results on large cores and egalitarian solutions, a coalitionally monotonic
core selection exists on four-player exact games. The
Housman--Clark mechanism is provably unavailable on exact games of any size,
common-core intersection of bumped exact games can nevertheless fail (a
four-arm example whose family remains feasible), and by a compactness
argument existence for each $n$ is \emph{equivalent} to feasibility of all
finite bump families, which we conjecture for $n\ge5$.
\end{abstract}

\medskip
\noindent\textbf{Keywords:} cooperative games; core; coalitional
monotonicity; nucleolus; SD-prenucleolus; exact games.


\section{Introduction}

A transferable-utility (TU) cooperative game assigns to each coalition the
worth it can guarantee on its own; a single-valued solution recommends one
division of $v(N)$. Two desiderata have organized much of the literature on
such solutions. \emph{Core selection}: the recommendation should be stable
against coalitional deviation whenever stability is possible.
\emph{Coalitional monotonicity} \citep{Young85}: if exactly one
coalition's worth rises, no member of that coalition should receive less.

The two desiderata are incompatible. \citet{Young85} proved that for
$|N|\ge 5$ no core selection is coalitionally monotonic; \citet{HousmanClark98} extended the impossibility to $|N|\ge 4$ and showed it
is tight ($n=3$ admits an infinite family of monotonic core selections). The
nucleolus fails coalitional monotonicity even on four-player convex games
\citep{Hokari2000}, as do the per-capita nucleolus \citep{Arin2013}
and the core-center \citep{MirasCalvo2021}. Against this background,
\citet{ArinKatsev2011,ArinKatsev2014,ArinKatsev2016,ArinKatsev2018}
introduced the \emph{surplus-distributor prenucleolus}
($\SD$-prenucleolus),%
\footnote{The literature uses both ``SD-nucleolus'' (e.g.\ the survey of \citealp{Inarra2020}) and ``SD-prenucleolus'' (the titles of
\citealp{ArinKatsev2011,ArinKatsev2014}) for the Arin--Katsev construction; we
follow the latter. See Section~\ref{sec:sd-def}.}
a lexicographic core selection that \emph{is} coalitionally monotonic on
convex games and on monotonic games with veto players. On the domain of
convex games alone, egalitarian solutions in the tradition of \citet{DuttaRay89} are
also coalitionally monotonic core selections \citep{Dietzenbacher2021}; what
distinguishes the SD-prenucleolus is that it selects from the core on
\emph{all} balanced games while retaining monotonicity on both the convex
and the monotonic-veto domains. How far that good behavior extends beyond
these two classes was, to our knowledge, open.

This paper answers that question and, in doing so, isolates what we believe
is the right open problem in this area. Our contributions:

\begin{enumerate}[leftmargin=2em]
\item \textbf{Minimal impossibility (Section~\ref{sec:minimal}).}
Three-game impossibility families (a base and two bumps) are already
present in the literature: the textbook proof of
\citet[Theorem~8.2.3]{PelegSudholter2003} has exactly this shape, and the
contradiction in the construction of \citet{Arin2013} likewise uses only a
base and two of its bumps. What appears not to have been recorded is
optimality: we prove that two-game families are always feasible, so three
is \emph{exactly} the minimal size of an impossibility certificate, and we
give a self-contained three-game certificate whose infeasibility admits a
six-line proof.

\item \textbf{The SD boundary (Sections~\ref{sec:sd-def}--\ref{sec:sd-results}).}
The coalitional monotonicity of the SD-prenucleolus does not extend one step
beyond convexity: it fails on \emph{exact} games at $n=4$ (convex
$\subsetneq$ exact $\subsetneq$ totally balanced $\subsetneq$ balanced), with
a minimized integer witness. Three companion results complete the map: the
SD-prekernel is multi-valued on exact games (containing a positive-dimensional
core segment, so the convex-domain uniqueness theorem of \citet{ArinKatsev2018}
is tight); the SD-prenucleolus is \emph{not} population monotonic even on
convex games (witness at the minimal size $n=3$, with a short hand proof);
and on permutation games the SD-prenucleolus fails pairwise monotonicity on
the same witness as the classical nucleolus, with a larger violation.
Certificates for the computational theorems appear in
Appendix~\ref{app:certs}.

\item \textbf{The exact-games question (Section~\ref{sec:escape}).}
The Young/Housman--Clark impossibility does not descend to exact games with
at most four players: combining the large-core result of \citet{Biswas99} with Dietzenbacher's monotonic egalitarian core
selection on PES-stable games \citep{Dietzenbacher2021} yields a
coalitionally monotonic core selection on $\mathcal{E}_4$
(Proposition~\ref{prop:n4}). For general $n$ the question is open, and we
show the known impossibility mechanisms cannot decide it: the singleton-core
pinning mechanism cannot exist on exact games (Lemma~\ref{lem:singleton});
face bounds of groups are tight and invariant along monotonicity chains of
proper bumps (Theorem~\ref{thm:pin}); overlapping two-arm families always
admit a common core point (Lemma~\ref{lem:vescape}); yet common-core
intersection can fail at $m=n=4$ arms (Proposition~\ref{prop:cx}) while the
associated family remains feasible---so family feasibility, not core
intersection, is the right invariant. Exact bump edges contain whole exact
one-parameter arcs (Lemma~\ref{lem:exact-arcs}), and every minimal infeasible
family must be a connected acyclic labeled graph with a sparse Helly--Farkas
witness (Proposition~\ref{prop:minimal-structure}). A compactness argument
(Proposition~\ref{prop:compact}) shows that for each $n$, existence of a
coalitionally monotonic core selection on exact games is \emph{equivalent}
to the feasibility of all finite bump families, which we conjecture for
$n\ge5$ (Conjecture~\ref{conj:existence}). The procedural egalitarian route
cannot establish this conjecture on all exact games:
\citet[][Example~5]{Mauri2017} already gives an exact five-player game that
is not PES-stable.

\item \textbf{Market games (Section~\ref{sec:markets}).}
The nucleolus is \emph{not} pairwise monotonic on permutation games: raising
a single entry of the underlying value matrix can strictly lower the payoff
of an agent named by that entry (smallest witness found: $n=4$, integer
matrix of total mass 20). In contrast, on assignment games---the two-sided
special case---we found no violation for the nucleolus or the
SD-prenucleolus in an extensive float-screened search that includes one
fully exact component, an exhaustive certified sweep of a $2\times2$ grid. For the nucleolus this independently
supports a theorem announced in the 2008 OTKA project report of Solymosi,
Forg\'o and Pint\'er \citep{SFP2008}, whose proof we have not been able to
access; for the SD-prenucleolus the statement appears new and we pose it as a
conjecture. The permutation witness shows any proof must use two-sidedness
itself.
\end{enumerate}

All positive computational claims in this paper (witnesses, solution
values, class memberships) are certified in exact rational arithmetic;
lexicographic solutions are certified by Kohlberg-type balancedness criteria
and, for headline witnesses, re-derived by an algorithmically independent
exact solver. Negative search outcomes are float-screened, as detailed in
Section~\ref{sec:methods}; Appendix~\ref{app:certs} lists human-checkable
certificates, reproducible with the supplement's standalone verifier.

\section{Preliminaries}\label{sec:prelim}

A \emph{TU game} is a pair $(N,v)$ with $N=\{1,\dots,n\}$ and
$v:2^N\to\R$, $v(\emptyset)=0$. For $x\in\R^N$ and $S\subseteq N$ write
$x(S)=\sum_{i\in S}x_i$. The \emph{core} is
$\core(v)=\{x:\ x(N)=v(N),\ x(S)\ge v(S)\ \forall S\}$; $v$ is
\emph{balanced} if $\core(v)\neq\emptyset$, \emph{totally balanced} if every
subgame is balanced, \emph{exact} \citep{Schmeidler72} if for every $S$ there is
$x\in\core(v)$ with $x(S)=v(S)$, and \emph{convex} if
$v(S\cup T)+v(S\cap T)\ge v(S)+v(T)$ for all $S,T$. It is standard that
convex $\Rightarrow$ exact $\Rightarrow$ totally balanced $\Rightarrow$
balanced, all inclusions strict for $n\ge 4$ (respectively $n\ge 3$ for the
last).

A collection $\mathcal{B}$ of nonempty proper coalitions is \emph{balanced}
if there are strictly positive weights $(\lambda_S)_{S\in\mathcal{B}}$ with
$\sum_{S\in\mathcal{B}}\lambda_S\mathbf{1}_S=\mathbf{1}_N$. A
\emph{core selection} on a class $\mathcal{C}$ of balanced games is a map
$\theta$ with $\theta(v)\in\core(v)$ for all $v\in\mathcal{C}$.

\begin{definition}[Coalitional monotonicity \citep{Young85,HousmanClark98}]
\label{def:cmon}
A single-valued solution $\theta$ is \emph{coalitionally monotonic} on
$\mathcal{C}$ if for all $u,v\in\mathcal{C}$ with $u(T)\le v(T)$ for some
coalition $T$ and $u(S)=v(S)$ for all $S\neq T$, we have
$\theta_i(u)\le\theta_i(v)$ for all $i\in T$.
\end{definition}

Aggregate monotonicity is the special case $T=N$. \emph{Population
monotonicity} on a class closed under subgames requires
$\theta_i(v|_S)\le\theta_i(v|_T)$ for all $i\in S\subseteq T$; it suffices to
check $|T\setminus S|=1$.

Throughout, a \emph{bump} of $v$ at $(S,d)$, $d>0$, with
$\emptyset\neq S\subseteq N$, is the game $v+d\cdot e_S$ agreeing with $v$
everywhere except $(v+d\cdot e_S)(S)=v(S)+d$; the bump is \emph{proper} if
$S\neq N$. A \emph{(bump) family} is a finite set of games
together with bump relations between them; a family is \emph{feasible} if
one can choose a core element for each game satisfying all the monotonicity
inequalities of Definition~\ref{def:cmon} along the bump relations. The
existence of a monotonic core selection on a class immediately implies every
family inside the class is feasible; Proposition~\ref{prop:compact} states
the converse.

We will use repeatedly the following elementary consequence of exactness.

\begin{lemma}\label{lem:facebound}
Let $v$ be exact and $T\subseteq N$ nonempty. Then
$\max_{x\in\core(v)}x(T)=v(N)-v(N\setminus T)$. Moreover, for any partition
$\{T_1,\dots,T_m\}$ of $N$ with $m\ge 2$,
$\sum_k\bigl(v(N)-v(N\setminus T_k)\bigr)\ge v(N)$; for $m=1$ the same
inequality holds trivially with equality.
\end{lemma}

\begin{proof}
For $x\in\core(v)$, $x(T)=v(N)-x(N\setminus T)$, so maximizing $x(T)$ is
minimizing $x(N\setminus T)$; by exactness the minimum equals
$v(N\setminus T)$. For the second claim with $m\ge2$, the collection
$\{N\setminus T_k\}_k$ is balanced with weights $1/(m-1)$, so balancedness
of $v$ gives $\sum_k v(N\setminus T_k)\le(m-1)\,v(N)$, which rearranges to
the claim.
\end{proof}

\begin{lemma}[exact bump arcs]\label{lem:exact-arcs}
Let $v$ be exact and let $S$ be nonempty. If $S\subsetneq N$ and
$v'=v+d e_S$ is exact for some $d>0$, then $v+t e_S$ is exact for every
$0\leq t\leq d$. If $S=N$, then $v+t e_N$ is exact for every $t\geq0$.
\end{lemma}

\begin{proof}
First let $S\subsetneq N$ and write
$C_t=\core(v)\cap\{x:x(S)\geq v(S)+t\}$. This is the core of
$v+t e_S$ and contains $\core(v')$. Exactness of $v$ supplies a point of
its core with $x(S)=v(S)$, while any point of $\core(v')$ has
$x(S)\geq v(S)+d$. Convexity therefore gives a point of $C_t$ tight at
$S$. For $T\neq S$, the function
\[
 f_T(t)=\min_{x\in C_t}x(T)
\]
is nondecreasing, is at least $v(T)$, and equals $v(T)$ at both endpoints
by exactness of $v$ and $v'$. Hence it is identically $v(T)$.

Now let $S=N$. For each proper $T$, choose
$x^T\in\core(v)$ with $x^T(T)=v(T)$ and a player $j\notin T$.
Then $x^T+t e_j$ belongs to the core of $v+t e_N$ and remains tight at
$T$. Efficiency handles $N$ itself.
\end{proof}

Thus every bump edge inside the exact class represents a whole exact
one-parameter arc, not merely a relation between two isolated games. This
permits local or parametric selection arguments to be integrated along the
same edges that define the finite-family feasibility problem.

\section{Minimal impossibility on balanced games}\label{sec:minimal}

\citet{HousmanClark98} prove that no core selection on four-player
balanced games is coalitionally monotonic, using a family of five games (a
base and four bumps). Certificates with three games are in fact already
present in the literature: the textbook proof of
\citet[Theorem~8.2.3]{PelegSudholter2003} uses precisely a four-player base
game with two single-coalition increases, and the contradiction in the
construction of \citet{Arin2013} also needs only a base and two of its
bumps. The observation that appears to be new here is the matching lower
bound: three games cannot be improved upon.

\begin{lemma}\label{lem:twogames}
Every two-game family $\{v,\,v+d\cdot e_S\}$ of balanced games is feasible.
\end{lemma}

\begin{proof}
If $S\neq N$, then $\core(v+d\cdot e_S)\subseteq\core(v)$, since the bump
only tightens the constraint at $S$; pick any $x\in\core(v+d\cdot e_S)$ and
assign it to both games. If $S=N$, pick $x\in\core(v)$ and assign
$x+\tfrac{d}{n}\mathbf{1}$ to the bumped game: it is efficient for it,
satisfies every core constraint (each $x(T)$ only increases), and dominates
$x$ coordinatewise. In both cases all monotonicity inequalities hold.
\end{proof}

\begin{theorem}\label{thm:threegame}
Let $N=\{1,2,3,4\}$ and define the balanced game $v$ by
\[
v(\{1,3\})=v(\{2,3\})=v(\{1,4\})=v(\{2,4\})=v(\{3,4\})=3,
\]
\[
v(\{1,2,3\})=v(\{1,2,4\})=v(\{1,3,4\})=1,\quad v(\{2,3,4\})=5,\quad v(N)=6,
\]
all other coalitions $0$. Let $v_A=v+3\cdot e_{\{1,2,3\}}$ and
$v_B=v+1\cdot e_{\{2,3,4\}}$; all three games are balanced, as witnessed by
$(1,1,2,2)\in\core(v)\cap\core(v_A)$ and $(0,0,3,3)\in\core(v_B)$. Then the
family $\{v,v_A,v_B\}$ is infeasible: no choice of core elements
$x\in\core(v)$, $y\in\core(v_A)$, $z\in\core(v_B)$ satisfies
$x_i\le y_i$ $(i\in\{1,2,3\})$ and $x_i\le z_i$ $(i\in\{2,3,4\})$.
Consequently no coalitionally monotonic core selection exists on any class
of games containing these three games; by Lemma~\ref{lem:twogames}, three
is the minimal possible size of such a certificate.
\end{theorem}

\begin{proof}
In $\core(v_A)$: adding $y(\{1,2,3\})\ge 4$, $y(\{1,4\})\ge 3$ and
$y(\{2,4\})\ge 3$ gives $2y_1+2y_2+y_3+2y_4\ge 10$, i.e.\
$2\,v_A(N)-y_3\ge 10$, so $y_3\le 2$. In $\core(v_B)$: adding
$z(\{1,3\})\ge 3$, $z(\{1,4\})\ge 3$ and $z(\{2,3,4\})\ge 6$ gives
$2z_1+z_2+2z_3+2z_4\ge 12$, so $z_2\le 0$. Monotonicity then forces
$x_3\le y_3\le 2$ and $x_2\le z_2\le 0$, whence
$x(\{2,3\})\le 2<3=v(\{2,3\})$, contradicting $x\in\core(v)$. Minimality
follows from Lemma~\ref{lem:twogames}.
\end{proof}

\begin{remark}
The proof is a rearrangement of an exact Farkas certificate with eleven
multipliers (nine core and monotonicity constraints with weight $1$; the two
efficiency identities with weight $2$), found by linear programming over
the family (Section~\ref{sec:methods}; the multiplier table is in the
supplement). None of the three games is exact, consistent with the
obstructions of Section~\ref{sec:escape}.
\end{remark}

\section{The SD-prenucleolus}\label{sec:sd-def}

\subsection{Definition}

Fix $(N,v)$ and an efficient $x$ ($x(N)=v(N)$). The SD-prenucleolus of
Arin--Katsev \citep{ArinKatsev2011,ArinKatsev2014} is defined through a
recursive \emph{satisfaction} vector, as follows. Let
$\mathcal{C}=2^N\setminus\{\emptyset,N\}$. Maintain a set
$H\subseteq\mathcal{C}$ of coalitions whose satisfactions $F(T;x)$ have been
fixed. For $S\notin H$ let
$H(S)=\bigcup\{T\in H:\ T\subseteq S\}$ (a set of \emph{players}) and, for
$i\in H(S)$,
\[
f(i,S)\;=\;\min\{F(T;x):\ T\in H,\ i\in T\subseteq S\}.
\]
The current candidate satisfaction of $S$ is
\[
F_H(S)=
\begin{cases}
\dfrac{x(S)-v(S)-\sum_{i\in H(S)}f(i,S)}{|S|-|H(S)|}, & H(S)\neq S,\\[1.4em]
x(S)-v(S)-\sum_{i\in S}f(i,S)+\max_{i\in S}f(i,S), & H(S)=S.
\end{cases}
\]
Starting from $H=\emptyset$ (all candidates are then per-capita
satisfactions $(x(S)-v(S))/|S|$), repeatedly fix \emph{all} coalitions
attaining the minimal current candidate at that value and update; the
procedure terminates with the satisfaction vector $F(\cdot\,;x)$. The
\emph{SD-prenucleolus} $\SD(N,v)$ is the efficient $x$ whose non-decreasingly
sorted satisfaction vector is lexicographically maximal. It is single-valued,
continuous, covariant, a core selection on balanced games, strongly aggregate
monotonic, and coalitionally monotonic on convex games and on monotonic veto
games \citep{ArinKatsev2011,ArinKatsev2014,ArinKatsev2016,ArinKatsev2018}.

The solution admits a characterization in the style of
\citet{Kohlberg71} \citep{ArinKatsev2011,ArinKatsev2018}: an efficient $x$ equals $\SD(N,v)$ if
and only if for every threshold $\alpha\in\R$ the lower level set
$\{S\in\mathcal{C}:\ F(S;x)\le\alpha\}$ is empty or balanced. We use this as
the certification instrument throughout.

\subsection{Implementation and validation}\label{sec:sd-impl}

We compute $\SD(N,v)$ by a Maschler-type successive linear-programming
scheme (at each stage the candidate satisfactions of unfixed coalitions are
affine in $x$, because fixed satisfactions are constants), with every output
certified by (i) an exact evaluation of the recursion above at the candidate
point and (ii) the Kohlberg-type criterion, with balancedness of each level
set established by an exact linear program producing strictly positive
weights. As additional regression checks (not proofs), headline witnesses
were probed for lexicographic optimality along hundreds of exact rational
perturbation directions.

The implementation reproduces exactly: the worked example of
\citet{ArinKatsev2018}; the prenucleolus on two-player games; the serial rule
on random monotonic veto games (Theorem 14.3 of \citealp{ArinKatsev2018}, using
the explicit serial-rule formula of \citealp{ArinFeltkamp2012}); the
minimal-overlap division on the bankruptcy game
$(E;c)=(200;100,200,300)$, namely $(100/3,\,250/3,\,250/3)$ (Theorem 14.4 of
\citealp{ArinKatsev2018}), which differs from the Talmud/nucleolus division
$(50,75,75)$; strict strong aggregate monotonicity on $150$ random games; and
coalitional monotonicity along $8$ random convex bumps. On the five-player
benchmark game of Arin--I\~narra reported in \citet{Inarra2020} our classical
nucleolus solver reproduces the published value
$(16/9,16/9,16/9,4/3,4/3)$ exactly.

\section{The boundary of the SD-prenucleolus}\label{sec:sd-results}

\subsection{Coalitional monotonicity fails on exact games}

\begin{theorem}\label{thm:sdexact}
The SD-prenucleolus is not coalitionally monotonic on the class of exact
games, already for $n=4$. Witness: let $v$ be
\[
v(\{2,3\})=1,\ v(\{1,2,3\})=1,\ v(\{2,4\})=13,\ v(\{1,2,4\})=15,
\]
\[
v(\{1,3,4\})=14,\ v(\{2,3,4\})=13,\ v(N)=30,
\]
all other coalitions $0$, and let $v'=v+1\cdot e_{\{1,3,4\}}$. Both games are
exact (hence totally balanced; neither is convex). Then
\[
\SD(v)=\Bigl(\tfrac{53}{12},\ \tfrac{19}{4},\ \tfrac{15}{4},\ \tfrac{205}{12}\Bigr),
\qquad
\SD(v')=\Bigl(\tfrac{35}{8},\ \tfrac{33}{8},\ \tfrac{33}{8},\ \tfrac{139}{8}\Bigr),
\]
so player $1\in\{1,3,4\}$ loses $1/24$ when the worth of a coalition
containing her rises.
\end{theorem}

\begin{proof}[Proof (computer-assisted)]
Exactness of each game is certified by $2^n-2$ exact linear programs
exhibiting, for every coalition $S$, a core element with $x(S)=v(S)$. That
the displayed points are the SD-prenucleoli follows from the Kohlberg-type
criterion of Section~\ref{sec:sd-def}: the full satisfaction level structure
of both points and strictly positive balancing weights for the lower level
sets are listed in Appendix~\ref{app:certs} and can be checked by hand. The
payoff comparison is exact rational arithmetic.
\end{proof}

The witness was found by search and greedily minimized (total mass $87$,
seven nonzero values); both games are moreover superadditive, so the failure
holds on exact (hence totally balanced) superadditive games. Since
Arin--Katsev prove monotonicity
on convex games, the boundary along
convex $\subsetneq$ exact $\subsetneq$ totally balanced $\subsetneq$ balanced
is \emph{convexity itself}.

\begin{corollary}
Exactness does not imply SD-relevance, the domain hypothesis of the general
monotonicity theorem of \citet{ArinKatsev2018}.\footnote{We refer to
\citet{ArinKatsev2018} for the definition of SD-relevant games; here we use
only the fact that their theorem establishes coalitional monotonicity of the
SD-prenucleolus on that class, which by Theorem~\ref{thm:sdexact} therefore
cannot contain both witness games.}
\end{corollary}

\subsection{The SD-prekernel is multi-valued on exact games}

\citet{ArinKatsev2016b,ArinKatsev2018} define the SD-prekernel
through bilateral complaints
$s_{ij}(x)=\min\{F(S;x):\ i\in S,\ j\notin S\}$: it is the set of efficient
$x$ with $s_{ij}(x)=s_{ji}(x)$ for all $i\neq j$. It contains the
SD-prenucleolus, and on convex games it equals it (a singleton).

\begin{theorem}\label{thm:sdpk}
The SD-prekernel can be multi-valued---indeed positive-dimensional---on
exact games. For the exact, totally balanced, non-convex game
\[
\begin{aligned}
&v(\{1\})=5,\ v(\{2\})=6,\ v(\{3\})=2,\ v(\{4\})=5,\\
&v(\{1,2\})=20,\ v(\{1,3\})=10,\ v(\{2,3\})=8,\ v(\{1,4\})=12,\
 v(\{2,4\})=14,\ v(\{3,4\})=12,\\
&v(\{1,2,3\})=22,\ v(\{1,2,4\})=25,\ v(\{1,3,4\})=17,\ v(\{2,3,4\})=18,\quad
 v(N)=32,
\end{aligned}
\]
the SD-prekernel contains the segment
$\SD(v)+t\,(1,-1,-1,1)$, $t\in[-1/2,1/2]$, where
$\SD(v)=(19/2,\,21/2,\,9/2,\,15/2)$; every point of the segment lies in the
core and equalizes all six bilateral complaints exactly.
\end{theorem}

\begin{proof}
Write $x^t=\SD(v)+t\,(1,-1,-1,1)$ for $t\in[-1/2,1/2]$. Core membership of
$x^t$ reduces to that of the endpoints $(9,11,5,7)$ and $(10,10,4,8)$, which
is a direct check against the fifteen constraints. Evaluating the recursion
of Section~\ref{sec:sd-def} at $x^t$ gives, for every $t\in[-1/2,1/2]$,
\[
F_{\{1,2\}}=F_{\{3,4\}}=0,\qquad F_{\{1,3\}}=F_{\{2,4\}}=2,
\]
\[
F_{\{3\}}=F_{\{1,2,3\}}=\tfrac52-t,\qquad
F_{\{4\}}=F_{\{1,4\}}=F_{\{1,2,4\}}=\tfrac52+t,
\]
\[
F_{\{2\}}=F_{\{2,3\}}=F_{\{2,3,4\}}=\tfrac92-t,\qquad
F_{\{1\}}=F_{\{1,3,4\}}=\tfrac92+t
\]
The displayed functions are proved by running the recursion symbolically,
stage by stage, on each open half-interval; the fixing order is constant on
each. For $t\in(-\tfrac12,0)$ the stages are
\[
\{12,34\};\ \{13,24\};\ \{4,14,124\};\ \{3,123\};\ \{1,134\};\
\{2,23,234\}
\]
at levels $0,\,2,\,\tfrac52+t,\,\tfrac52-t,\,\tfrac92+t,\,\tfrac92-t$
respectively, and for $t\in(0,\tfrac12)$ they are
\[
\{12,34\};\ \{13,24\};\ \{3,123\};\ \{4,14,124\};\ \{2,23,234\};\
\{1,134\}
\]
at levels $0,\,2,\,\tfrac52-t,\,\tfrac52+t,\,\tfrac92-t,\,\tfrac92+t$. At
each stage, direct substitution of $x^t$ into the two-case formula of
Section~\ref{sec:sd-def}, using the already-fixed values, yields exactly the
displayed candidates and confirms that the stated minimizers are fixed next
(the required inequalities, e.g.\ $2\le\tfrac52\mp t$ and
$\tfrac52+t<\tfrac52-t$ for $t<0$, hold strictly on the open
half-intervals). The boundary values $t\in\{-\tfrac12,0,\tfrac12\}$ are
verified separately by direct exact computation (supplement), where some
stages merge but the resulting satisfaction values still obey the displayed
formulas. Reading off the bilateral complaints
$s_{ij}=\min\{F_S:\ i\in S\not\ni j\}$: for $|t|\le\tfrac12$,
$s_{12}=s_{21}=2$ (attained at $\{1,3\}$ resp.\ $\{2,4\}$),
$s_{34}=s_{43}=2$, and all other ordered pairs have $s_{ij}=0$ (attained at
$\{1,2\}$ or $\{3,4\}$). All six equalities $s_{ij}=s_{ji}$ therefore hold
along the segment. That the midpoint is $\SD(v)$ is certified by the
Kohlberg-type criterion (Appendix~\ref{app:certs}).
\end{proof}

Thus the uniqueness theorem for convex games is tight, with the same
boundary as Theorem~\ref{thm:sdexact}. (Perturbed control points off the
segment fail the equalities, confirming the checker discriminates.)

\subsection{Population monotonicity fails on convex games}

\begin{theorem}\label{thm:sdpm}
The SD-prenucleolus is not population monotonic on convex games, already at
$n=3$, which is minimal: for $n\le 2$ every core selection is population
monotonic, since individual rationality gives
$\theta_i(v|_{\{i\}})=v(\{i\})\le\theta_i(v|_T)$ for any $T\ni i$. For the
convex, exact game
$v(\{1,2\})=v(\{2,3\})=1$, $v(N)=2$ (all else $0$):
\[
\SD\bigl(v|_{\{2,3\}}\bigr)=\bigl(\tfrac12,\tfrac12\bigr),\qquad
\SD(v)=\bigl(\tfrac13,\ \tfrac43,\ \tfrac13\bigr),
\]
so player $3$ falls from $1/2$ to $1/3$ when player $1$ joins.
\end{theorem}

\begin{proof}
Convexity is immediate. The two-player subgame on $\{2,3\}$ has
$v(\{2\})=v(\{3\})=0$, $v(\{2,3\})=1$; the SD-prenucleolus coincides with
the prenucleolus on two-player games, giving the equal split $(1/2,1/2)$.

For the full game we verify $x=(1/3,4/3,1/3)$ by the Kohlberg-type
criterion. The recursion finishes in two stages (the five coalitions at the
minimal level are fixed first, then $\{2\}$): the per-capita satisfactions
are $F(\{1\})=F(\{3\})=1/3$;
$F(\{1,2\})=F(\{2,3\})=(x(S)-1)/2=1/3$; $F(\{1,3\})=x(\{1,3\})/2=1/3$;
$F(\{2\})=4/3$. The lower level sets are: for $\alpha\in[1/3,4/3)$,
$\mathcal{B}_1=\{\{1\},\{3\},\{1,2\},\{2,3\},\{1,3\}\}$, balanced with
strictly positive weights
$\lambda_{\{1,2\}}=\lambda_{\{2,3\}}=\tfrac12$,
$\lambda_{\{1,3\}}=\tfrac13$, $\lambda_{\{1\}}=\lambda_{\{3\}}=\tfrac16$
(player $1$: $\tfrac12+\tfrac13+\tfrac16=1$; player $2$:
$\tfrac12+\tfrac12=1$; player $3$: $\tfrac13+\tfrac12+\tfrac16=1$); for
$\alpha\ge 4/3$, the set of all proper coalitions, which is balanced. Hence
every lower level set is empty or balanced and $x=\SD(v)$. Finally
$1/2>1/3$.
\end{proof}

In exploratory float-screened sampling of convex games (integer-valued
supermodular games generated as nonnegative integer combinations of
unanimity games), violations appeared in roughly half of the instances at
$n=4$ and more at $n=5$; we report these frequencies as indicative only,
since the exploratory counting runs were not preserved as seeded scripts.
The phenomenon is in any case common under such sampling, not exceptional. The SD-prenucleolus therefore
separates the two monotonicities on its home class: it is coalitionally
monotonic there \citep{ArinKatsev2016} but not population monotonic---in
contrast with the egalitarian solution of \citet{DuttaRay89}, which is
population monotonic on convex games.

\section{Market games: pairwise monotonicity}\label{sec:markets}

An \emph{assignment game} \citep{NunezRafels2002} is generated by a
nonnegative matrix $A=(a_{ij})$ over two disjoint sides $M$ (rows) and $M'$
(columns): the player set is $M\cup M'$ and $v(S)$ is the maximum weight of
a matching between $S\cap M$ and $S\cap M'$ with weights $a_{ij}$. A
\emph{permutation game} \citep{TijsEtAl84} is generated by a square matrix
$K=(k_{ij})$ over a single player set:
$v(S)=\max_\pi\sum_{i\in S}k_{i\pi(i)}$, the maximum over permutations
$\pi$ of $S$. Assignment games embed into permutation games
\citep{CurielTijs86}; both classes are totally balanced.

\emph{Pairwise monotonicity} \citep{NunezRafels2002} is
the natural monotonicity axiom for such games: if one entry of the matrix
rises, neither of the two agents named by that entry may lose. The
buyers-optimal, sellers-optimal and $\tau$-value rules satisfy it on
assignment games \citep{NunezRafels2002,vdBNunezRobles2021}; \citet{Miquel2009} constructs a bespoke pairwise-monotonic core selection for
permutation games.%
\footnote{We could not access the body of \citet{Miquel2009}; its abstract
states a positive result for a merged-$\tau$-value rule and makes no claim
about the nucleolus. Whether an explicit nucleolus counterexample appears in
its introduction is a caveat for the novelty (not the correctness) of
Theorem~\ref{thm:permutation}.}
For the nucleolus on assignment games, the 2008 OTKA project report of
Solymosi, Forg\'o and Pint\'er \citep{SFP2008} announces: ``We proved that
the nucleolus is pairwise-monotonic on assignment games''; the project's
2012 final report attributes the result to an unpublished manuscript,
``Pairwise monotonicity of the nucleolus in assignment games'' by
T.~Solymosi, M.~Brugueras and T.~E.~S.~Raghavan (2012), with a simplified
proof in a revised version.\footnote{OTKA final report,
\url{https://real.mtak.hu/12261/1/72856_ZJ1.pdf}.} We have not been able to
access either manuscript; our computations below independently support the
statement.

\begin{theorem}\label{thm:permutation}
The nucleolus is not pairwise monotonic on permutation games. Let $K$ be the
value matrix (rows: owners; columns: recipients)
\[
K=\begin{pmatrix}0&7&0&0\\ 1&0&0&1\\ 1&0&0&0\\ 0&9&1&0\end{pmatrix},
\]
and raise $K_{43}$ from $1$ to $2$ (which raises only
$v(\{3,4\})$ and $v(\{1,3,4\})$, each by one, and fixes $v(N)=11$). The
nucleolus moves from
$(\tfrac12,\ \tfrac{33}{4},\ \tfrac12,\ \tfrac74)$ to
$(\tfrac23,\ \tfrac{23}{3},\ \tfrac13,\ \tfrac73)$: player $3$, a member of
both improved coalitions, falls from $1/2$ to $1/3$. Both games are totally
balanced and both nucleoli lie in the core. Moreover the SD-prenucleolus
fails on the same instance, falling from $1/2$ to $1/4$ for the same player.
\end{theorem}

\begin{proof}[Proof (computer-assisted)]
Both nucleolus claims are certified by Kohlberg's criterion in exact
arithmetic; Appendix~\ref{app:certs} lists the excess level structure and
balancing weights at the top level for both games, and the machine-readable
supplement contains the full level data. The SD claims are certified as in
Section~\ref{sec:sd-impl}.
\end{proof}

The witness is greedily minimized (integer matrix, total mass $20$); a
float-screened $n=3$ sweep of $36{,}000$ bumps found no candidate
violation, so $n=4$ is the smallest size at which we found the phenomenon. In this instance the improving pair's
\emph{joint} payoff rises ($9/4\to 8/3$ for the nucleolus), so the witness
does not violate Sasaki's weak (joint) pairwise monotonicity \citep{Sasaki95};
the bump's surplus is redistributed \emph{within} the improving pair.

On assignment games, by contrast, an extensive search---float-screened
random and tie-targeted $3\times3$, $4\times3$, $4\times4$ markets and
hill-climbing on the violation margin, together with one fully exact
component, an exhaustive Kohlberg-certified sweep of all $2\times2$ markets
with entries in $\{0,\dots,4\}$ (2{,}500 bumps, zero violations, zero
certification failures)---found no violation for either solution. For the nucleolus this is consistent with \citep{SFP2008}; for the
SD-prenucleolus we know of no prior statement:

\begin{conjecture}\label{conj:assign}
The SD-prenucleolus is pairwise monotonic on assignment games.
\end{conjecture}

\begin{conjecture}\label{conj:mcst}
The nucleolus of minimum-cost spanning tree games is cost monotonic in the
sense of \citet{DuttaKar2004} (raising one edge's cost never
lowers an incident player's cost share).
\end{conjecture}

Conjecture~\ref{conj:mcst} is supported by $1{,}450$ float-screened tests
(with exact certification awaiting any candidate violation; none arose) at
$n\in\{4,5,6\}$ including large bumps in the core-moving regime of
\citet{TanCaoZou2025}; the nucleolus is known to be monotonic on the
fixed-tree subclass \citep{MPR2010}, and Dutta--Kar's violating rule is
Bird's, not the nucleolus. Theorem~\ref{thm:permutation} shows any proof of
the assignment-game statements must use two-sidedness: the property fails
one class up.

\section{Do exact games escape the impossibility?}\label{sec:escape}

The certificate games of Theorem~\ref{thm:threegame}, like those of
\citet{HousmanClark98}, are not exact. This section first records that for
$n\le4$ the existence question on exact games is settled affirmatively by
known results (Proposition~\ref{prop:n4}), then shows that the pinning and
two-arm mechanisms underlying the known impossibility proofs are
unavailable on exact games of any size, delimits the techniques by an
explicit counterexample to a natural strengthening, and reduces the
existence question to finite-family feasibility.

\subsection{The pinning mechanism is unavailable}

The Housman--Clark proof pins players via bumps whose cores collapse to
singletons.

\begin{lemma}\label{lem:singleton}
An exact game whose core is a singleton is additive. Consequently no exact
game family can pin a player through a singleton core except trivially.
\end{lemma}

\begin{proof}
If $\core(v)=\{x\}$, exactness gives, for each $S$, a core element attaining
$x(S)=v(S)$; but the core element is unique, so $x(S)=v(S)$ for every $S$,
i.e.\ $v$ is additive with $v(S)=\sum_{i\in S}x_i$.
\end{proof}

More generally, pinning arguments transmit \emph{face bounds}
$\max_{x\in\core}x(T)$ along monotonicity chains. On exact games these
bounds cannot contract under proper bumps:

\begin{theorem}[Pin invariance under proper bumps]\label{thm:pin}
Let all games below be exact and let all bumps be proper. Fix a nonempty
group $T\subseteq N$. Along any monotonicity chain that constrains
$T$---i.e.\ any sequence of proper bumps at coalitions $S\supseteq T$---the
face bound of Lemma~\ref{lem:facebound} satisfies
$\max_{x\in\core}x(T)=v(N)-v(N\setminus T)$ in every game of the chain, with
both $v(N)$ and $v(N\setminus T)$ unaltered by every bump in the chain. In
particular, since face bounds over any partition sum to at least $v(N)$
(Lemma~\ref{lem:facebound}), no contradiction with efficiency can be derived
from partition-aggregated face bounds, at any chain depth and for groups of
any arity; we do not claim the analogue for arbitrary balanced collections
of groups. (A bump at $S=N$ raises every face bound by $d$ and is likewise
harmless.)
\end{theorem}

\begin{proof}
A proper bump at $S$ constrains $T$ only if $T\subseteq S$; since
$T\cap(N\setminus T)=\emptyset$ and $S\neq N$, no such $S$ equals
$N\setminus T$ or $N$, so neither $v(N\setminus T)$ nor $v(N)$---and with
them the face bound, by Lemma~\ref{lem:facebound} applied in each (exact)
game---changes along the chain.
\end{proof}

\subsection{Two-arm families: positive results and their limits}

\begin{lemma}[Overlapping arms]\label{lem:vescape}
Let $b$ be exact and let $t_A=b+d_A e_A$, $t_B=b+d_B e_B$ be exact, where
$A\neq B$ are nonempty proper coalitions with $A\cap B\neq\emptyset$. Then
$M=b+d_Ae_A+d_Be_B$ is balanced, and any $x^*\in\core(t_A)$ maximizing
$x(B)$ lies in
$\core(M)\subseteq\core(b)\cap\core(t_A)\cap\core(t_B)$; assigning $x^*$ to
all three games shows the family $\{b,t_A,t_B\}$ is feasible.
\end{lemma}

\begin{proof}
By exactness of $t_A$ at $N\setminus B$ and $A\not\subseteq N\setminus B$
(as $A\cap B\neq\emptyset$),
\[
\max_{x\in\core(t_A)}x(B)
= t_A(N)-t_A(N\setminus B)
= b(N)-b(N\setminus B)
= \max_{x\in\core(b)}x(B)\;\ge\; b(B)+d_B,
\]
the last inequality because $\core(t_B)\neq\emptyset$ forces some core
element of $b$ with $x(B)\ge b(B)+d_B$. The maximizer $x^*$ then satisfies
every constraint of $M$.
\end{proof}

\begin{lemma}[Singleton arms are transparent]\label{lem:singleton-arm}
Let $b$ and $t_i=b+d e_{\{i\}}$ be balanced.  For every
$x\in\core(b)$ there is $y\in\core(t_i)$ with $y_i\geq x_i$.
Consequently, in a three-game star
$\{b,t_A,t_B\}$, if one arm is a singleton then the family is feasible
(without any exactness assumption).
\end{lemma}

\begin{proof}
Put $u=\max\{z_i:z\in\core(b)\}$.  Balancedness of $t_i$ means that
some core point of $b$ satisfies $z_i\geq b(\{i\})+d$, so
$u\geq b(\{i\})+d$.  A maximizer $y$ of $z_i$ over $\core(b)$ therefore
belongs to $\core(t_i)$ and satisfies $y_i=u\geq x_i$ for every
$x\in\core(b)$.

For the consequence, suppose $A=\{i\}$.  Choose any
$x\in\core(t_B)\subseteq\core(b)$ and assign it to both $b$ and $t_B$.
The first statement supplies $y\in\core(t_A)$ with $y_i\geq x_i$, which
completes the family witness.
\end{proof}

\begin{corollary}[complete reduction for two-arm stars]
\label{cor:two-arm-reduction}
Let $b,t_A,t_B$ be exact single-bump games.  The star family
$\{b,t_A,t_B\}$ is feasible whenever $A\cap B\neq\emptyset$ or
$\min\{|A|,|B|\}=1$.  Hence every such family on at most three players is
feasible.  Any unresolved two-arm obstruction must have
\[
 A\cap B=\emptyset,\qquad |A|\geq2,\qquad |B|\geq2,
\]
so it begins with a two-pair partition at $n=4$.
\end{corollary}

\begin{proof}
Use Lemma~\ref{lem:vescape} in the overlapping case and
Lemma~\ref{lem:singleton-arm} in the singleton case.  Two disjoint nonempty
coalitions on at most three players cannot both have size at least two.
\end{proof}

\begin{proposition}[Disjoint arms: exactness chaining]\label{prop:disjoint}
Let $b$ be exact and $t_A=b+d_Ae_A$, $t_B=b+d_Be_B$ exact, where $A,B$ are
nonempty proper coalitions with $A\cap B=\emptyset$ and
$R:=N\setminus(A\cup B)\neq\emptyset$. Then
\[
b(N)\ \ge\ b(A)+d_A+b(B)+d_B+b(R),
\]
i.e.\ the partition Farkas obstruction $\{A,B,R\}$ to the balancedness of
$M=b+d_Ae_A+d_Be_B$ vanishes. (This is one obstruction among many; the
proposition does not by itself establish balancedness of $M$ or feasibility
of the family.)
\end{proposition}

\begin{proof}
\emph{Step 1.} Every $x\in\core(t_B)$ has
$x(N\setminus A)=x(B)+x(R)\ge b(B)+d_B+b(R)$; since $t_B$ is exact at
$N\setminus A$ and $t_B(N\setminus A)=b(N\setminus A)$ (here
$B\subsetneq N\setminus A$ because $R\neq\emptyset$, so the bump at $B$ does
not change the worth of $N\setminus A$), taking the minimum gives
$b(N\setminus A)\ge b(B)+d_B+b(R)$.
\emph{Step 2.} By exactness of $t_A$ at $N\setminus A$ there is
$x^*\in\core(t_A)$ with $x^*(N\setminus A)=b(N\setminus A)$, hence
$x^*(A)=b(N)-b(N\setminus A)$; but $x^*(A)\ge b(A)+d_A$, so
$b(N)\ge b(A)+d_A+b(N\setminus A)$. Chaining the two steps gives the claim.
\end{proof}

These techniques have a genuine limit, which the following example (due to a
reviewer of an earlier draft) makes precise: common-core intersection of
bumped games---the conclusion of Lemma~\ref{lem:vescape}---can fail when
$m=n=4$ arms are bumped, even though the associated monotonicity family
remains feasible.

\begin{proposition}[Core intersection can fail at $m=n=4$; feasibility
survives]\label{prop:cx}
Let $n=4$, $b(N)=5$ and $b(S)=0$ for all proper $S$. For $k\in N$ let
$t_k=b+4\,e_{N\setminus\{k\}}$. Then $b$ and each $t_k$ are exact, and
$\core(t_k)=\{x\ge 0:\ x(N)=5,\ x_k\le 1\}$, so
$\bigcap_{k}\core(t_k)=\emptyset$ (four inequalities $x_k\le1$ would force
$x(N)\le4$); equivalently, the all-arm game
$b+\sum_k 4\,e_{N\setminus\{k\}}$ is unbalanced. Nevertheless the star
family $\{b,t_1,\dots,t_4\}$ is feasible: take $\theta(b)=(2,1,1,1)$,
$\theta(t_1)=(1,\tfrac43,\tfrac43,\tfrac43)$, and for $k\neq 1$ the point of
$\core(t_k)$ obtained from $(2,1,1,1)$ by moving one unit from player $k$ to
player $1$ (e.g.\ $\theta(t_2)=(3,0,1,1)$), which dominates $\theta(b)$ on
$N\setminus\{k\}$.
\end{proposition}

\begin{proof}
Exactness of $b$ is clear (all mass can be placed on any single player).
For $t_k$: the core is as displayed. For a proper $S\ni k$, the point with
$x_k=0$ and mass $5$ on a player outside $S$ attains $x(S)=0=t_k(S)$. For a
proper $S\not\ni k$ with $S\neq N\setminus\{k\}$, there is a player
$j\notin S\cup\{k\}$; the point with $x_k=1$, $x_j=4$ attains
$x(S)=0=t_k(S)$. Finally the minimum of $x(N\setminus\{k\})=5-x_k$ over the
core is $4=t_k(N\setminus\{k\})$. So $t_k$ is exact. Emptiness of the
intersection and unbalancedness (the collection $\{N\setminus\{k\}\}_k$ with
weights $1/3$ gives $\tfrac{16}3>5$) are immediate. Feasibility: each
displayed $\theta(t_k)$ lies in $\core(t_k)$ and dominates
$\theta(b)=(2,1,1,1)$ on every coordinate of $N\setminus\{k\}$.
\end{proof}

Proposition~\ref{prop:cx} corrects a conjecture we advertised in an earlier
draft (that the all-arm game of exact single bumps is always balanced): the
common-point method cannot decide the exact-games question, because family
feasibility is strictly weaker than core intersection. Computationally,
across more than $10^6$ family-feasibility linear programs built from over
$7{,}500$ exact four-player bases (two- and three-arm stars and depth-two
lattices), floating-point screening flagged \emph{no infeasible family of
exact games}. The pipeline certifies any flagged infeasibility by an exact
Farkas certificate (validated on the Housman--Clark family and on
Theorem~\ref{thm:threegame}); the negative outcome itself is
float-screen-level, and float screening can in principle produce false
negatives near degeneracy. In light of Proposition~\ref{prop:n4} below, the
four-player census outcome is in fact implied by the existence of a
monotonic selection on $\mathcal{E}_4$, so we present it as validation of
the feasibility pipeline rather than as evidence at the open frontier,
which begins at $n=5$.

\subsection{Equivalence, and the state of the question}

\begin{proposition}\label{prop:compact}
Fix $n$ and let $\mathcal{E}_n$ be the class of exact $n$-player games. A
coalitionally monotonic core selection exists on $\mathcal{E}_n$ if and only
if every finite bump family within $\mathcal{E}_n$ is feasible.
\end{proposition}

\begin{proof}
Necessity is immediate. For sufficiency, consider the product
$P=\prod_{v\in\mathcal{E}_n}\core(v)$, compact in the product topology by
Tychonoff's theorem, since each core is a nonempty compact subset of
$\R^n$. For each bump pair $(u,v)$ in $\mathcal{E}_n$ let
$C_{u,v}\subseteq P$ be the set of selections satisfying the (finitely many,
closed) monotonicity inequalities of Definition~\ref{def:cmon} for that
pair; $C_{u,v}$ is closed and depends on two coordinates only. A finite
intersection of such sets involves finitely many games; feasibility of the
corresponding finite family provides a point (extending arbitrarily
elsewhere). By the finite intersection property the full intersection is
nonempty, and any of its points is a coalitionally monotonic core selection.
\end{proof}

\begin{proposition}[structure of a minimal finite obstruction]
\label{prop:minimal-structure}
Suppose a finite bump family on $q$ games is infeasible and is
inclusion-minimal with respect to its games and bump relations. Label each
directed bump edge by its bumped coalition. Then:
\begin{enumerate}[label=(\roman*),leftmargin=2em]
\item the underlying undirected graph is connected and the directed graph is
acyclic;
\item for every edge $u\to w$ labeled $S$, some $i\in S$ has no alternative
directed path from $u$ to $w$ all of whose edge labels contain $i$;
\item infeasibility is witnessed by at most $q(n-1)+1$ core or monotonicity
halfspaces, after the $q$ efficiency equalities are imposed. For rational
games, the corresponding sparse Farkas certificate may be chosen rational.
\end{enumerate}
\end{proposition}

\begin{proof}
If the undirected graph were disconnected, feasible selections for its
components could be combined, contrary to minimal infeasibility. Along every
directed edge exactly one coalition worth strictly increases, so the sum of
all coalition worths is a strictly increasing potential; directed cycles are
impossible.

For (ii), suppose instead that for every $i\in S$ there were such an
alternative path. Transitivity of the monotonicity inequalities along that
path would imply $\theta_i(u)\leq\theta_i(w)$. Thus every inequality attached
to the edge $u\to w$ would follow from the remaining edges, contradicting
minimality.

Finally restrict the $qn$ allocation variables to the product of the $q$
efficiency hyperplanes, an affine space of dimension $q(n-1)$. Family
feasibility is the intersection there of finitely many core and monotonicity
halfspaces. Helly's theorem gives an empty subintersection of at most
$q(n-1)+1$ halfspaces. With rational coefficients, the usual rational form
of Farkas' lemma supplies rational multipliers on such a subfamily.
\end{proof}

Proposition~\ref{prop:minimal-structure} turns the search for a refuting exact
family into a sparse labeled-DAG problem. In particular, disconnected stars,
directed cycles, transitively redundant bump labels, and certificates whose
support exceeds the Helly bound need not be searched when mining minimal
obstructions.

For $n\le 4$ the existence question is in fact already settled in the
affirmative by combining two known results:

Recall that the core of $v$ is \emph{large} \citep{Sharkey82} if for
every $y\in\R^N$ with $y(S)\ge v(S)$ for all $S$ there is
$x\in\core(v)$ with $x\le y$, and call a game \emph{PES-stable} if the
procedural egalitarian solution of \citet{Dietzenbacher2021} selects a
core allocation of it.

\begin{proposition}[\citealp{Dietzenbacher2021}, with \citealp{Biswas99}]
\label{prop:n4}
A coalitionally monotonic core selection exists on the class
$\mathcal{E}_4$ of exact four-player games (hence also on $\mathcal{E}_n$,
$n\le 4$).
\end{proposition}

\begin{proof}
Every exact game with at most four players has a large core
\citep{Biswas99}. Large-core games are PES-stable, and on PES-stable games
the procedural egalitarian solution is a single-valued core selection
satisfying coalitional monotonicity \citep{Dietzenbacher2021}; the
combination is observed in \citet{Dietzenbacher2021}. (As a consistency
check, we verified by exact vertex enumeration that all eight exact
four-player games directly certified in this paper have large cores.)
\end{proof}

Proposition~\ref{prop:compact} frames both sides of the general ledger: the
impossibility on balanced games \emph{is} one infeasible finite family
(Theorem~\ref{thm:threegame}); the results above are partial progress
toward feasibility of all finite exact families for general $n$. The
solutions studied in this paper fail coalitional monotonicity on exact
games individually---the nucleolus (already on convex games
\citep{Hokari2000}), the per-capita nucleolus \citep{Arin2013}, the
core-center \citep{MirasCalvo2021} (and, in our tests, a coordinate-extreme
variant, in $136$ of $1{,}059$ exact bumps), and the SD-prenucleolus
(Theorem~\ref{thm:sdexact})---while Proposition~\ref{prop:n4} shows that
for $n\le4$ a different, egalitarian-type formula succeeds. For $n\ge5$ no
selection is known to work. Dietzenbacher's particular procedural
egalitarian solution cannot supply one on all exact games:
\citet[][Example~5]{Mauri2017} applies the same maximum-average procedure to
an exact five-player game, obtains the zero vector, and observes that it lies
outside the core. The recent follow-up of \citet{DietzDogan2025} develops
egalitarianism under monotonicity constraints further but does not settle the
existence of a different selection on exact games.

\begin{conjecture}\label{conj:existence}
For every fixed $n\ge5$, a coalitionally monotonic core selection exists on
the class $\mathcal{E}_n$ of exact $n$-player games.
\end{conjecture}

By Proposition~\ref{prop:compact}, Conjecture~\ref{conj:existence} is
equivalent to feasibility of all finite exact families for the given $n$.
The pinning-based and partition-aggregated obstruction mechanisms are
unavailable at any $n$ (Lemma~\ref{lem:singleton}, Theorem~\ref{thm:pin});
what is missing is a
feasibility proof for arbitrary finite families---or a refuting family.
Proposition~\ref{prop:cx} shows that one natural route to refutation, empty
common intersection of bumped cores, is by itself insufficient. An
argument for arbitrary finite families---or a new finite obstruction---is
therefore still required.

\section{Verification methodology}\label{sec:methods}

All results were produced by a search-then-certify discipline. Screening
uses floating-point successive linear programs; every reported object is
then reconstructed in exact rational arithmetic and certified independently
of the search path: nucleolus and SD points by Kohlberg-type
balanced-collection criteria evaluated by an exact rational simplex method
(written for this project after we found that a current computer-algebra
system's exact linear-programming routine returns infeasible points as
optima on three-variable instances; our solver is fuzz-tested against a
standard floating-point solver with exact feasibility assertions), class
memberships (balancedness, total balancedness, exactness, convexity) by
exact linear programs, family infeasibility by exact Farkas certificates
re-verified under constraint reordering, and headline witnesses additionally
re-derived by an algorithmically independent exact successive-LP solver.
Negative search outcomes (no violation found, no infeasible family found)
are reported at the float-screen level: exact certification in those
campaigns is triggered by candidate positives, and none arose; the sole
fully exact negative sweep is the exhaustive $2\times2$ assignment grid.
Known values from the literature (Talmud division, serial rule,
minimal-overlap rule, a published five-player nucleolus, Hokari-type
failures at their known minimal sizes, and the Housman--Clark family itself)
are reproduced as anchors. The discipline caught nine false or overstated
candidate results before any claim was circulated, including two
floating-point artifacts near the exactness boundary and the overreaching
core-intersection conjecture corrected in Proposition~\ref{prop:cx}; we
accordingly treat direct exact LP oracles as authoritative over derived
closed forms. The computational search, the verification tooling, and the
drafting of this paper made substantial use of large-language-model
assistance; every mathematical claim was verified by the certification
pipeline described here, and all text was reviewed by the author.
Computations used Python 3.10 with exact arithmetic via the
standard \texttt{fractions} module and floating-point screening via
NumPy~2.2/SciPy~1.15 (HiGHS). A machine-readable supplement
(\texttt{certificates.json}) contains, for every computational theorem, the
full game tables, per-coalition exactness witnesses, complete
satisfaction/excess tables, and strictly positive balancing weights for
every cumulative level set, together with a one-command verifier
(\texttt{verify\_paper.py}) that re-checks all of it in exact arithmetic
using only the Python standard library (it recomputes every satisfaction
and excess table from scratch, checks that every stored level set equals
the recomputed one, and validates the supplied balancing weights by pure
rational arithmetic), and an environment manifest; together these form
Online Resource~1.

\section{Open problems}\label{sec:open}

\begin{enumerate}[leftmargin=2em]
\item Prove (or refute) feasibility of every finite bump family of exact
$n$-player games for $n\ge5$; by Proposition~\ref{prop:compact} this
settles Conjecture~\ref{conj:existence}. Proposition~\ref{prop:cx} shows
the common-core-point method cannot decide it, and the procedural egalitarian
route is closed by the exact counterexample of \citet[][Example~5]{Mauri2017}.
\item Recover the unpublished Solymosi--Brugueras--Raghavan proof
(see Section~\ref{sec:markets}) that the nucleolus is pairwise monotonic on
assignment games, and prove
Conjecture~\ref{conj:assign} for the SD-prenucleolus.
Theorem~\ref{thm:permutation} shows two-sidedness must be used.
\item Prove Conjecture~\ref{conj:mcst} (Dutta--Kar cost monotonicity of the
nucleolus on MCST games), known on the fixed-tree subclass \citep{MPR2010}.
\item Characterize SD-relevance \citep{ArinKatsev2018} within the exact
games; Theorem~\ref{thm:sdexact} shows exactness is insufficient.
\item Determine whether the witness of Theorem~\ref{thm:permutation} is
pre-empted by an example in the introduction of \citet{Miquel2009}.
\end{enumerate}

\subsection*{Statements and Declarations}
\small
\textbf{Funding:} none.
\textbf{Data and code availability:} a machine-readable certificate
supplement accompanies this paper as Online Resource~1
(\texttt{certificates.json}; \texttt{verify\_paper.py};
\texttt{environment.txt}); the verifier is self-contained and uses only
the Python standard library.
\textbf{Use of AI assistance:} see Section~\ref{sec:methods}.
\textbf{Competing interests:} none.
\normalsize

\appendix

\section{Certificates}\label{app:certs}

All certificates below are exact rational data; balancing weights are
strictly positive multipliers $\lambda$ with
$\sum_S\lambda_S\mathbf{1}_S=\mathbf{1}_N$. This appendix displays the key
tables; the \emph{complete} certificate package---per-coalition exactness
witnesses for the eight games claimed exact by direct certification (the
convex game of Theorem~\ref{thm:sdpm} is exact by convexity), full
satisfaction and excess tables, and balancing weights for \emph{every}
cumulative level set of every solution point claimed (including the SD
points of the permutation games and the two-player subgame of
Theorem~\ref{thm:sdpm})---is in the machine-readable supplement
(\texttt{certificates.json}), re-verifiable in exact arithmetic, using only
the Python standard library, with the single command
\texttt{python3 verify\_paper.py}.

\subsection*{Theorem~\ref{thm:sdexact}}

Satisfaction levels of $\SD(v)$ (ascending; coalitions at each level):
\[
\begin{array}{ll}
\alpha=15/4:& \{3\},\ \{2,3\},\ \{1,2,4\},\ \{1,3,4\}\\
\alpha=53/12:& \{1\},\ \{1,3\},\ \{1,2,3\},\ \{2,4\}\\
\alpha=19/4:& \{2\},\ \{1,2\}\\
\alpha=61/12:& \{2,3,4\}\\
\alpha=205/12:& \{4\},\ \{1,4\},\ \{3,4\}
\end{array}
\]
The minimal level set is balanced with weights
$\lambda_{\{3\}}=\tfrac13$, $\lambda_{\{2,3\}}=\tfrac13$,
$\lambda_{\{1,2,4\}}=\tfrac23$, $\lambda_{\{1,3,4\}}=\tfrac13$. For
$\SD(v')$ the levels are
\[
\begin{array}{ll}
\alpha=29/8:& \{2,3\},\ \{1,2,4\},\ \{1,3,4\}\\
\alpha=33/8:& \{2\},\ \{3\}\\
\alpha=35/8:& \{1\},\ \{1,2\},\ \{1,3\},\ \{1,2,3\},\ \{2,4\}\\
\alpha=43/8:& \{2,3,4\}\\
\alpha=139/8:& \{4\},\ \{1,4\},\ \{3,4\}
\end{array}
\]
with minimal level set balanced by
$\lambda_{\{2,3\}}=\lambda_{\{1,2,4\}}=\lambda_{\{1,3,4\}}=\tfrac12$.
Cumulative lower level sets are balanced at every threshold (weights in the
supplement), certifying both points by the criterion of
Section~\ref{sec:sd-def}.

\subsection*{Theorem~\ref{thm:sdpk}}

At every sampled point of the segment ($t=-\tfrac12,-\tfrac13,-\tfrac14,
0,\tfrac14,\tfrac13,\tfrac12$, all in the supplement) the bilateral
complaints are
\[
s_{12}=s_{21}=2,\qquad s_{34}=s_{43}=2,\qquad
s_{ij}=0\ \text{for all other ordered pairs},
\]
and the two half-interval stage orders displayed in the proof of
Theorem~\ref{thm:sdpk} extend these equalities to the whole segment.

\subsection*{Theorem~\ref{thm:permutation}}

For the base game $K$: the nucleolus
$(\tfrac12,\tfrac{33}4,\tfrac12,\tfrac74)$ has maximal excess $0$, attained
exactly by the collection $\{\{1,3\},\{2,4\}\}$, balanced with weights
$\lambda=1,1$ (a partition). For the bumped game $K'$: the nucleolus
$(\tfrac23,\tfrac{23}3,\tfrac13,\tfrac73)$ likewise has maximal excess $0$
attained by the partition $\{\{1,3\},\{2,4\}\}$. Kohlberg's criterion
requires every cumulative upper excess set to be balanced; the supplement
lists all eight (resp.\ seven) thresholds with strictly positive weights for
both games, together with the complete SD certificates for the SD subclaim.

\begin{thebibliography}{99}

\bibitem[Arin(2013)]{Arin2013}
Arin J (2013) Monotonic core solutions: beyond Young's theorem.
Int J Game Theory 42:325--337.
\url{https://doi.org/10.1007/s00182-013-0368-8}

\bibitem[Arin and Feltkamp(2012)]{ArinFeltkamp2012}
Arin J, Feltkamp V (2012) Coalitional games: monotonicity and core.
Eur J Oper Res 216:208--213.
\url{https://doi.org/10.1016/j.ejor.2011.07.025}

\bibitem[Arin and Katsev(2011)]{ArinKatsev2011}
Arin J, Katsev I (2011) The SD-prenucleolus for TU games.
IKERLANAK Working Paper IL~56/11, University of the Basque Country

\bibitem[Arin and Katsev(2014)]{ArinKatsev2014}
Arin J, Katsev I (2014) The SD-prenucleolus for TU games.
Math Methods Oper Res 80:307--327.
\url{https://doi.org/10.1007/s00186-014-0482-9}

\bibitem[Arin and Katsev(2016a)]{ArinKatsev2016}
Arin J, Katsev I (2016a) A monotonic core solution for convex TU games.
Int J Game Theory 45:1013--1029.
\url{https://doi.org/10.1007/s00182-015-0500-z}

\bibitem[Arin and Katsev(2016b)]{ArinKatsev2016b}
Arin J, Katsev I (2016b) The SD-prekernel for TU games.
IKERLANAK Working Paper IL~96/16, University of the Basque Country

\bibitem[Arin and Katsev(2018)]{ArinKatsev2018}
Arin J, Katsev I (2018) The SD-prenucleolus for TU-games: coalitional
monotonicity and core stability. In: Mueller D, Trost R (eds) Game Theory
in Management Accounting. Springer, Cham, pp 301--321.
\url{https://doi.org/10.1007/978-3-319-61603-2_14}

\bibitem[Biswas et al.(1999)]{Biswas99}
Biswas AK, Parthasarathy T, Potters JAM, Voorneveld M (1999) Large cores
and exactness. Games Econ Behav 28:1--12.
\url{https://doi.org/10.1006/game.1998.0686}

\bibitem[Curiel and Tijs(1986)]{CurielTijs86}
Curiel I, Tijs S (1986) Assignment games and permutation games.
Methods Oper Res 54:323--334

\bibitem[Dietzenbacher(2021)]{Dietzenbacher2021}
Dietzenbacher B (2021) Monotonicity and egalitarianism.
Games Econ Behav 127:194--205.
\url{https://doi.org/10.1016/j.geb.2021.03.006}

\bibitem[Dietzenbacher and Do\u{g}an(2025)]{DietzDogan2025}
Dietzenbacher B, Do\u{g}an E (2025) Population monotonicity and
egalitarianism. Int J Game Theory 54:42.
\url{https://doi.org/10.1007/s00182-025-00966-y}

\bibitem[Dutta and Kar(2004)]{DuttaKar2004}
Dutta B, Kar A (2004) Cost monotonicity, consistency and minimum cost
spanning tree games. Games Econ Behav 48:223--248.
\url{https://doi.org/10.1016/j.geb.2003.09.008}

\bibitem[Dutta and Ray(1989)]{DuttaRay89}
Dutta B, Ray D (1989) A concept of egalitarianism under participation
constraints. Econometrica 57:615--635.
\url{https://doi.org/10.2307/1911055}

\bibitem[Hokari(2000)]{Hokari2000}
Hokari T (2000) The nucleolus is not aggregate-monotonic on the domain of
convex games. Int J Game Theory 29:133--137.
\url{https://doi.org/10.1007/s001829900030}

\bibitem[Housman and Clark(1998)]{HousmanClark98}
Housman D, Clark L (1998) Note: core and monotonic allocation methods.
Int J Game Theory 27:611--616.
\url{https://doi.org/10.1007/s001820050093}

\bibitem[I\~narra et al.(2020)]{Inarra2020}
I\~narra E, Serrano R, Shimomura KI (2020) The nucleolus, the kernel, and
the bargaining set: an update. Rev \'Econ 71:225--266.
\url{https://doi.org/10.3917/reco.712.0225}

\bibitem[Kohlberg(1971)]{Kohlberg71}
Kohlberg E (1971) On the nucleolus of a characteristic function game.
SIAM J Appl Math 20:62--66.
\url{https://doi.org/10.1137/0120009}

\bibitem[Maschler et al.(2010)]{MPR2010}
Maschler M, Potters J, Reijnierse H (2010) The nucleolus of a standard
tree game revisited: a study of its monotonicity and computational
properties. Int J Game Theory 39:89--104.
\url{https://doi.org/10.1007/s00182-009-0176-3}

\bibitem[Mauri Masdeu(2017)]{Mauri2017}
Mauri Masdeu L (2017) Essays on egalitarianism-based solution concepts for
cooperative TU-games. Doctoral dissertation, Universitat Rovira i Virgili.
\url{http://hdl.handle.net/10803/460900}

\bibitem[Miquel(2009)]{Miquel2009}
Miquel S (2009) A pairwise-monotonic core selection for permutation games.
Math Methods Oper Res 70:465--475.
\url{https://doi.org/10.1007/s00186-008-0281-2}

\bibitem[Mir\'as Calvo et al.(2021)]{MirasCalvo2021}
Mir\'as Calvo MA, Quinteiro Sandomingo C, S\'anchez-Rodr\'iguez E (2021)
Considerations on the aggregate monotonicity of the nucleolus and the
core-center. Math Methods Oper Res 93:291--325.
\url{https://doi.org/10.1007/s00186-020-00733-7}

\bibitem[Peleg and Sudh\"olter(2003)]{PelegSudholter2003}
Peleg B, Sudh\"olter P (2003) Introduction to the Theory of Cooperative
Games. Kluwer Academic Publishers, Boston.
\url{https://doi.org/10.1007/978-1-4615-0308-8}

\bibitem[N\'u\~nez and Rafels(2003)]{NunezRafels2002}
N\'u\~nez M, Rafels C (2003) The assignment game: the $\tau$-value.
Int J Game Theory 31:411--422.
\url{https://doi.org/10.1007/s001820300127}

\bibitem[Sasaki(1995)]{Sasaki95}
Sasaki H (1995) Consistency and monotonicity in assignment problems.
Int J Game Theory 24:373--397.
\url{https://doi.org/10.1007/BF01243039}

\bibitem[Schmeidler(1972)]{Schmeidler72}
Schmeidler D (1972) Cores of exact games, I.
J Math Anal Appl 40:214--225.
\url{https://doi.org/10.1016/0022-247X(72)90045-5}

\bibitem[Sharkey(1982)]{Sharkey82}
Sharkey WW (1982) Cooperative games with large cores.
Int J Game Theory 11:175--182.
\url{https://doi.org/10.1007/BF01755727}

\bibitem[Solymosi et al.(2008)]{SFP2008}
Solymosi T, Forg\'o F, Pint\'er M (2008) J\'at\'ekelm\'elet [Game theory].
OTKA project report. \url{https://real.mtak.hu/1330/}

\bibitem[Tan et al.(2025)]{TanCaoZou2025}
Tan Z, Cao Z, Zou Z (2025) Comparative statics of minimum-cost spanning
tree games. Games Econ Behav 151:162--182.
\url{https://doi.org/10.1016/j.geb.2025.03.005}

\bibitem[Tijs et al.(1984)]{TijsEtAl84}
Tijs SH, Parthasarathy T, Potters JAM, Rajendra Prasad V (1984)
Permutation games: another class of totally balanced games.
OR Spektrum 6:119--123.
\url{https://doi.org/10.1007/BF01721088}

\bibitem[van den Brink et al.(2021)]{vdBNunezRobles2021}
van den Brink R, N\'u\~nez M, Robles F (2021) Valuation monotonicity,
fairness and stability in assignment problems.
J Econ Theory 195:105277.
\url{https://doi.org/10.1016/j.jet.2021.105277}

\bibitem[Young(1985)]{Young85}
Young HP (1985) Monotonic solutions of cooperative games.
Int J Game Theory 14:65--72.
\url{https://doi.org/10.1007/BF01769885}

\end{thebibliography}

\end{document}
